\section{Introduction}\label{sec:introduction}

Machine-learning models can redistribute causal responsibility when their internal computation is perturbed. Ablating one component can change the contribution of others, producing self-repair, backup behavior, or more general conditional importance. For mechanistic interpretation, this creates a methodological problem: detecting that a component became more or less important does not identify why its contribution changed. A selected attention head may alter its task-relevant routing scores, its competing scores, or its value content; alternatively, essentially the same candidate computation may become more or less useful because the surrounding network has changed. Reliable mechanistic evidence therefore requires separating the detection of conditional importance from the explanation of its mechanism.

We introduce \textbf{FALI (Factorized Attention-Level Intervention)}, a hypothesis-driven methodology for resolving this ambiguity. FALI begins after a candidate attention head and a task-relevant target edge have been identified. It compares an intact state with a prespecified perturbed state and decomposes the resulting change in the candidate's \emph{output-grounded causal usefulness}. The candidate's attention computation is factorized into target scores $T$, rival scores $R$, other structurally valid scores $O$, value content $V$, and surrounding-network context $C$. Evaluating the resulting $2^5$ counterfactual states defines a causal game whose interactions are allocated with the standard Owen value. The fine score contributions reconcile exactly with a coarser SCORE/VALUE/CONTEXT decomposition.

The key distinction is between \emph{localizing conditional importance} and \emph{explaining it}. Activation and path patching identify behaviorally relevant components or paths; conditional co-ablation identifies components whose importance changes after a primary intervention; interaction-aware mediation analyses show that a component's measured effect can depend strongly on the rest of the model. FALI is downstream of these approaches: given a selected conditional change, it asks which architecture-aligned changes account for it under an explicit intervention model. This makes FALI a mechanism-resolution methodology rather than an automatic circuit-discovery procedure.

We first validate the intervention semantics in controlled systems. Deterministic synthetic attention models isolate each of $T,R,O,V,C$ and a known cross-factor interaction, while a small randomly initialized transformer checks score-to-pattern reconstruction, value-to-output reconstruction, endpoint recovery, intervention isolation, batch-specific score partitions, and consistency between the fine Owen game and its coarse quotient game. These tests establish that the implementation recovers the factors it manipulates; they do not imply that the factorization is unique or that mixed hybrid states are naturally occurring model states.

We then apply FALI to GPT-2 Small on indirect object identification (IOI). Corrected functional analysis separates attention adaptation from recruitment: after the canonical primary Name Movers are removed, L10H2 becomes more causally useful, whereas L10H10 increases target attention even more strongly while its causal usefulness decreases. FALI assigns both heads a target-positive/rival-negative score pattern, showing that even a favorable-looking score reorganization does not itself define recruitment. For positively recruited L10H2, the main fine-score contributions are $\phi_T>0$ and $\phi_R<0$ across the canonical confirmation cohort.

The decisive test is cross-setting transfer. We prospectively test the GPT-2 target-positive/rival-negative mechanism in Pythia-160M induction. Positive conditional recruitment appears again, but the fine score signs reverse: $\phi_T<0$ and $\phi_R>0$. The mechanism-transfer hypothesis therefore fails. The shared functional label---conditional recruitment---does not imply a shared fine-grained mechanism.

Our contributions are threefold. First, we define an output-grounded counterfactual estimand for the cross-state change in a selected attention head's causal usefulness. Second, we factor that change into score, value, and surrounding-context contributions, with target/rival/other resolution inside SCORE and coarse--fine consistency via standard Owen attribution. Third, we show empirically that two settings with positive conditional recruitment can exhibit opposite target/rival score mechanisms. The contribution is therefore a research methodology for discriminating mechanisms of conditional importance in learned neural networks, not a new theory of self-repair, a new localization method, or new Shapley/Owen mathematics.

An earlier public preprint from this project used the name FALI for a narrower \textbf{Frozen Attention Logit Intervention} analysis of attention-probability changes in GPT-2 IOI. The present paper substantially revises and supersedes that formulation: it changes the estimand to output-grounded causal usefulness, expands the intervention game to $T/R/O/V/C$, adds controlled implementation validation and coarse--fine reconciliation, broadens the GPT-2 confirmation, and adds a prospectively specified cross-setting mechanism-transfer test. The detailed relationship is stated in Related Work and in the MLJ Contribution Information Sheet.
